\section{Evaluation \& Agent Improvement Process}
\subsection{Rubric-Based Evaluation Criteria}
The evaluation framework was designed to measure both the quality of the generated database artifacts and the quality of the agent process used to produce them. This separation follows agent-evaluation guidance that recommends assessing not only final answers, but also reasoning traces, tool use, and intermediate actions \cite{deepeval-agent-eval,google-agent-eval,braintrust-agent-eval,openlayer-agent-eval}. Therefore, each task was evaluated with two independent scores: an \textbf{artifact quality score} and an \textbf{agent process quality score}. The two scores were intentionally not merged, because a task can produce a correct deliverable while still using an incomplete, inefficient, or unsafe workflow.

\paragraph{Artifact quality rubric.}
Artifact quality was scored from 0 to 5 using a weighted database-design rubric. A score of 0 means the criterion is missing or seriously wrong, while a score of 5 means the output is strong, complete, consistent, and well justified. The weighted artifact score is calculated as the sum of each criterion score multiplied by its weight.

\begin{table}[htbp]
\centering
\small
\begin{tabularx}{\textwidth}{p{0.28\textwidth}p{0.12\textwidth}X}
\toprule
\textbf{Criterion} & \textbf{Weight} & \textbf{Evaluation focus} \\
\midrule
Completeness of required outputs & 10\% & Required files exist and contain all mandatory sections for the active task. \\
Requirement coverage & 20\% & Source requirements and business rules are preserved, especially booking conflicts, unavailable-space restrictions, status constraints, and maintenance rules. \\
Database design correctness & 20\% & Entities, relationships, cardinalities, tables, keys, constraints, triggers, and indexes are technically reasonable and enforce the intended rules. \\
Traceability & 15\% & Requirements can be followed through business rules, entities, relationships, tables, constraints, sample data, and queries. \\
Assumptions and open questions & 10\% & Missing or ambiguous information is explicitly recorded instead of being silently invented. \\
Cross-file consistency & 10\% & Entity names, table names, attributes, statuses, and constraints remain consistent across outputs, registries, and design decisions. \\
SQL Server compatibility & 10\% & SQL artifacts use valid Microsoft SQL Server syntax and provide execution evidence where required. \\
Query usefulness & 5\% & Final queries answer realistic operational needs such as booking history, utilization, no-shows, maintenance, and unavailable spaces. \\
\bottomrule
\end{tabularx}
\caption{Artifact quality rubric used for generated database design outputs}
\label{tab:artifact-rubric}
\end{table}

\paragraph{Agent process rubric.}
Process quality was evaluated from the trajectory logs rather than from the final artifact alone. This matches trajectory-based evaluation for multi-step agents, where silent process failures can occur even when the final answer appears plausible \cite{google-agent-eval,openlayer-agent-eval}. After generating or revising an \texttt{outputs/0X-*} artifact, the agent records the plan, actual execution steps, files read and written, deviations, outcome, and self-detected fixes. The six process metrics are scored independently from 0 to 5 and then averaged into a process score.

\begin{table}[htbp]
\centering
\small
\begin{tabularx}{\textwidth}{p{0.27\textwidth}X}
\toprule
\textbf{Metric} & \textbf{Evaluation focus} \\
\midrule
PlanQuality & Whether the plan is logical, ordered, complete, and suited to the active pipeline task. \\
PlanAdherence & Whether actual execution followed the plan or justified deviations clearly. \\
StepEfficiency & Whether the agent avoided redundant reads, writes, retries, and unnecessary context loading. \\
TaskCompletion & Whether the required output exists and satisfies the task's success criteria. \\
ToolCorrectness & Whether the agent read and wrote the expected files while avoiding forbidden paths. \\
ArgumentCorrectness & Whether commands and file targets used correct paths, group suffixes, and read/write directions. \\
\bottomrule
\end{tabularx}
\caption{Agent process metrics scored from trajectory logs}
\label{tab:process-rubric}
\end{table}

\paragraph{Evaluation workflow.}
The recommended workflow is iterative: the agent first generates a task artifact, then writes a trajectory log, then the evaluator runs a task-level evaluation when a quality check is needed. The evaluation report returns an artifact score, a process score, supporting evidence, and the top improvement actions. Those results are then used to revise the output, template, command, or task-specific skill when needed. At larger milestones, a pipeline-level evaluation checks cross-task consistency. This reflects the practice of combining end-to-end artifact evaluation with component-level checks for tool choice, file access, arguments, and side effects \cite{braintrust-agent-eval,openlayer-agent-eval,dailydose-agent-metrics}.



\subsection{Agent Improvement Techniques}
\subsubsection{Task-Scoped Progressive Context Disclosure}
The first improvement technique was \textbf{task-scoped progressive context disclosure}. Instead of loading the entire repository, the agent followed the documented reading order and opened only the files relevant to the active pipeline stage. The process started with project rules and memory files, then moved to task-specific skills, templates, and the exact upstream artifacts required for the current deliverable.

This technique matched the seven-task pipeline. For Task~01, the authoritative input was \texttt{req/business-requirement.md}, so the agent avoided premature schema assumptions. For Tasks~02 and 03, the context expanded to the entity registry, ERD, business-rule analysis, and technical conventions. After schema freeze, Tasks~05 and 06 relied primarily on the locked schema registry, approved validation output, DDL, and task-specific SQL guidance.

This approach aligns with the progressive disclosure principle \cite{claudemem-progressive} and improved the agent in three ways: \textbf{First}, it reduced context pollution by excluding unrelated files. \textbf{Second}, it lowered cross-task drift via controlled inputs. \textbf{Third}, it simplified evaluation through verifiable logs.

\subsubsection{Command and Skill Metadata with YAML Front Matter}
The second improvement technique was the use of \textbf{YAML front matter} as a small machine-readable header before longer human-readable instructions. YAML front matter is a common pattern for storing file-level metadata at the beginning of Markdown-like documents \cite{jekyll-frontmatter}. In this project, it appeared in command files, skill files, and trajectory logs.

In \texttt{.opencode/commands/}, front matter helped identify command-level properties before the prompt body was read. In \texttt{.opencode/skills/}, it exposed \texttt{name} and \texttt{description} fields that helped the agent select task-specific behavior. In \texttt{logs/trajectory/}, it standardized run metadata such as \texttt{task}, \texttt{group}, \texttt{run\_at}, \texttt{status}, and \texttt{revision\_of}. OpenCode's command and skill documentation directly supports this style of command routing and skill discovery \cite{opencode-commands,opencode-skills}.

This improved the agent in three ways: \textbf{First}, it strengthened command routing through early command distinction. \textbf{Second}, it enhanced skill discovery without loading full task files. \textbf{Third}, it increased evaluation reliability via stable trajectory identifiers. To prevent metadata drift, schema facts and business rules were strictly maintained in authoritative registries rather than front matter.

\subsubsection{Registry-Grounded Source-of-Truth Control}
The third improvement technique was \textbf{registry-grounded source-of-truth control}. Stable database knowledge was externalized into persistent documentation: \texttt{docs/entity-registry.md} for conceptual entities and attributes, \texttt{docs/schema-registry.md} for relational tables and constraints, and \texttt{docs/design-decisions.md} for rationale and trade-offs.

This structure matched the project pipeline. Task~01 populated the entity registry from the business requirement. Task~02 refined relationships and cardinalities. Task~03 locked attributes and populated the schema registry. Task~04 validated both registries before schema freeze. After that point, Tasks~05--07 had to treat the registries as read-only control documents rather than silently modifying upstream design facts.

This improved the agent in three ways. \textbf{First}, it reduced context drift because later tasks did not need to rediscover the entity model from scratch. \textbf{Second}, it improved cross-file consistency because ERD, logical schema, DDL, sample data, and future query design could be checked against the same approved registries. \textbf{Third}, it introduced change control: if a later task needed to revise a design choice, the reason had to be recorded in \texttt{docs/design-decisions.md} rather than hidden inside a generated output.

This approach aligns with both \textbf{codified context infrastructure} and \textbf{source-of-truth synchronization} in agent systems. In our project, the registries acted as \textbf{domain-specific cold memory}, while the design-decision log served as a \textbf{governance record}. Instead of allowing duplicated local representations to drift, the agent must consult these canonical registries before producing or validating downstream artifacts. Together, this framework made the workflow more \textbf{auditable, reproducible, and resistant to hallucinated schema changes}.

\subsubsection{Requirement-to-Artifact Traceability Control}
The fourth improvement technique was \textbf{requirement-to-artifact traceability control}. Each important requirement needed a visible path through the pipeline: from the original business requirement, to numbered business rules, entities, relationships, logical tables, database constraints, DDL triggers, sample-data tests, and eventually business queries. This made it easier to verify that requirements were preserved instead of being dropped or reinterpreted across tasks.

In this project, traceability appeared as a distributed control mechanism rather than a single perfect matrix. Task~01 established numbered business rules and assumptions. The entity registry stored source references for conceptual entities and attributes. The schema registry mapped business rules to database-level enforcement mechanisms such as constraints, indexes, and triggers. Task~06 then strengthened downstream traceability by labeling expected-error cases with BR numbers and preserving SQL execution evidence.

This improved the agent in three ways. First, it enabled coverage checking by exposing missing requirement links. Second, it facilitated backward debugging by tracing design errors back to source rules. Third, it enhanced auditability by making the requirement's entire route inspectable. This is consistent with requirements traceability research, which treats trace links as a way to follow requirements forward and backward through development artifacts \cite{gotel-traceability}.




